% Chapter 1

\chapter{Introducción general}

\label{Chapter1}
\label{IntroGeneral}

En este capítulo se introduce el problema de la proyección de contrataciones a partir de información comercial histórica. Se presenta el rol del pronóstico de ventas en la planificación de una organización y se describen los desafíos propios de las ventas entre empresas. Luego se expone el contexto que motivó el trabajo y se formula el problema técnico. Finalmente, se revisan los principales enfoques propuestos en la literatura y se establecen los objetivos y el alcance del trabajo.

\section{Desafíos del pronóstico de ventas B2B}
\label{sec:desafiosPronosticoB2B}

La proyección de ventas constituye un insumo central para la planificación financiera de una organización. A partir de ella se estiman ingresos, se asignan recursos, se definen metas y se anticipan necesidades de capacidad. En entornos empresariales, esta tarea suele combinar información histórica con el conocimiento de los responsables comerciales. Los sistemas de pronóstico brindan una base cuantitativa, mientras que las personas incorporan información contextual que no siempre se encuentra representada en los datos. Sin embargo, los ajustes basados en juicio experto no mejoran necesariamente la precisión y pueden introducir sesgos, en particular cuando son frecuentes, pequeños o excesivamente optimistas \citep{Fildes2009}.

Los sistemas de gestión de relaciones con clientes o CRM (\textit{Customer Relationship Management}) registran de manera estructurada la actividad comercial. Para cada oportunidad pueden almacenar su monto esperado, la probabilidad estimada de cierre, la etapa dentro del proceso de venta, la fecha prevista de concreción y otras características relevantes. Además de facilitar el seguimiento operativo, estos registros permiten reconstruir cómo evolucionó el \textit{pipeline} comercial y ofrecen una fuente de datos para desarrollar pronósticos reproducibles.

En mercados de empresa a empresa o B2B (\textit{business to business}), el modelado presenta dificultades particulares. La cantidad de operaciones suele ser menor que en contextos de consumo masivo, los montos son heterogéneos y unas pocas oportunidades de gran tamaño pueden modificar de manera significativa el total de un período. También son frecuentes los cambios en las fechas esperadas de cierre y en las valoraciones subjetivas del equipo comercial. Se han propuesto modelos de aprendizaje automático para estimar la probabilidad de éxito de oportunidades B2B a partir de la evolución del \textit{pipeline} \citep{Yan2015}. Otros trabajos utilizaron información comercial anticipada para mejorar procesos manuales de pronóstico en casos reales de ventas B2B \citep{Rohaan2022}.

El trabajo aborda una variante diferente y complementaria. En lugar de estimar únicamente la probabilidad de éxito de cada oportunidad, busca proyectar el volumen agregado de contrataciones, entendido como el monto de ventas efectivamente comprometidas durante un horizonte futuro. La salida sirve como apoyo para la planificación financiera trimestral y permite una comparación directa con el pronóstico utilizado por la organización.

\section{Contexto y motivación}
\label{sec:contextoMotivacion}

OceanSound Partners es una firma de inversión en empresas de tecnología y servicios habilitados por tecnología que atienden mercados gubernamentales y empresariales \citep{OceanSoundWebsite}. Una de las empresas de su portafolio administra la actividad comercial mediante un CRM en el que registra la evolución de sus oportunidades de venta. La planificación financiera depende, en buena medida, de la proyección de las contrataciones futuras.

Antes de este trabajo, esas proyecciones se elaboraban a partir del pronóstico almacenado en el CRM y de la evaluación manual del equipo comercial. El proceso dependía del criterio de las personas involucradas, era difícil de reproducir y no contaba con un procedimiento común para medir su error histórico. Esta situación dificultaba anticipar desvíos en los ingresos, asignar recursos con suficiente previsión y distinguir si una corrección provenía de nueva información comercial o de una apreciación subjetiva. Como consecuencia, la organización tenía menos capacidad para detectar de manera temprana necesidades de contratación, ajustes presupuestarios y oportunidades que requerían atención.

El CRM conserva una fotografía semanal, denominada \textit{snapshot}, del estado completo del \textit{pipeline}. Una misma oportunidad aparece en sucesivos \textit{snapshots} con la información disponible en cada fecha. Esta estructura permite estudiar no solo el resultado final de una venta, sino también la evolución de las variables que podían observarse antes de que ese resultado ocurriera.

La motivación del trabajo surge de dos oportunidades. Desde la perspectiva organizacional, se busca disponer de un pronóstico cuantitativo, trazable y comparable con el método vigente. Desde la perspectiva académica, el problema permite integrar preparación y análisis de datos, aprendizaje automático, series temporales, validación de modelos y gestión de experimentos sobre un caso real.

\section{Descripción del problema}
\label{sec:descripcionProblema}

El conjunto de datos tiene estructura de panel longitudinal: cada fila representa el estado de una oportunidad en una fecha de observación y una misma oportunidad se repite a lo largo de múltiples \textit{snapshots}. Las contrataciones se determinan a partir de las oportunidades ganadas y de su fecha efectiva de cierre. Por esta razón, una agregación directa de todas las filas produciría duplicaciones y mezclaría estados conocidos en momentos diferentes.

El desafío metodológico principal consiste en respetar la disponibilidad temporal de la información. Para producir una proyección desde un instante $T$, las variables predictoras solo pueden utilizar datos que ya eran conocidos en ese instante. Incorporar el estado final de una oportunidad, su fecha de cierre efectiva o una transformación calculada con observaciones posteriores constituiría una fuga de información o \textit{data leakage}. Esta clase de error genera métricas artificialmente optimistas y un modelo que no puede reproducir su desempeño cuando se utiliza sobre períodos futuros.

La fuente combina así una gran cantidad de filas con una historia temporal relativamente corta. El problema requiere transformar los estados históricos del \textit{pipeline} en observaciones comparables, definir un monto de contrataciones sin contar una misma oportunidad más de una vez y relacionar cada observación con un período futuro. También exige considerar que los meses recientes pueden contener oportunidades cuyo resultado definitivo todavía no se encontraba registrado. El diseño adoptado para resolver estas dificultades se presenta en el capítulo~\ref{Chapter3}.

\section{Estado del arte}
\label{sec:estadoArte}

El pronóstico de ventas puede abordarse mediante métodos estadísticos, modelos de aprendizaje automático, redes neuronales y estimaciones basadas en juicio experto. No existe un enfoque universalmente superior: el resultado depende de la cantidad de datos, su frecuencia, el horizonte de decisión, la estabilidad del proceso y la disponibilidad de variables explicativas.

\subsection{Pronósticos estadísticos y basados en juicio experto}

Los modelos autorregresivos integrados de media móvil o ARIMA (\textit{Autoregressive Integrated Moving Average}) constituyen una referencia clásica para representar autocorrelación, tendencia y estacionalidad \citep{BoxJenkins2015}. Su formulación es interpretable y puede resultar efectiva cuando la estructura histórica de la serie permanece relativamente estable. Como limitación para este trabajo, un modelo puramente univariado no utiliza directamente el estado actual de las oportunidades comerciales.

Prophet propone un modelo descomponible con tendencia, estacionalidades y eventos, acompañado por parámetros que pueden ser ajustados por un analista \citep{TaylorLetham2018}. Fue concebido para facilitar el pronóstico de numerosas series de negocio y ofrece una referencia intermedia entre automatización e intervención humana. No obstante, tanto Prophet como ARIMA necesitan suficientes observaciones para identificar patrones temporales de manera confiable.

Los pronósticos basados en criterio comercial incorporan conocimiento que puede no estar codificado en los datos, pero su desempeño depende de cómo se realicen los ajustes. La evidencia muestra que algunos ajustes mejoran la precisión, mientras que otros la reducen y presentan sesgo optimista \citep{Fildes2009}. Por ese motivo, resulta necesario medir el valor agregado por la intervención humana y conservar el pronóstico del CRM como una referencia explícita.

\subsection{Aprendizaje automático aplicado al embudo comercial}

Los métodos basados en árboles permiten capturar relaciones no lineales e interacciones sin imponer una forma funcional única. Los bosques aleatorios o \textit{Random Forest} reducen la varianza mediante la combinación de árboles construidos sobre muestras y subconjuntos aleatorios de variables \citep{Breiman2001}. XGBoost agrega árboles de manera secuencial mediante \textit{gradient boosting} e incorpora regularización y optimizaciones para datos dispersos \citep{ChenGuestrin2016}. Ambos enfoques pueden utilizar rezagos históricos junto con medidas agregadas del \textit{pipeline}, aunque requieren control de complejidad cuando la cantidad de períodos es pequeña.

Las redes recurrentes LSTM (\textit{Long Short Term Memory}) fueron diseñadas para aprender dependencias prolongadas mediante mecanismos de memoria y compuertas \citep{HochreiterSchmidhuber1997}. Su flexibilidad permite representar dinámicas temporales complejas, pero implica una cantidad mayor de parámetros y una necesidad de datos que puede resultar desfavorable en series empresariales cortas. En este contexto, su inclusión constituye una hipótesis que debe evaluarse y no una garantía de mejora.

En el dominio B2B se utilizaron modelos para estimar la probabilidad de éxito de oportunidades a partir de información temporal del \textit{pipeline} \citep{Yan2015}. También se aplicó aprendizaje supervisado a información comercial anticipada para complementar procesos manuales de pronóstico \citep{Rohaan2022}. Estos antecedentes respaldan el uso de variables disponibles antes del cierre, aunque el objetivo del presente trabajo es el monto agregado de contrataciones y no la clasificación individual de oportunidades.

\subsection{Agregación temporal}

La agregación temporal transforma una serie de alta frecuencia en observaciones correspondientes a períodos más extensos. Esta operación atenúa variaciones de alta frecuencia y puede facilitar la identificación de tendencias y ciclos. La combinación de múltiples niveles de agregación permite mejorar la estimación de componentes estructurales de una serie \citep{Kourentzes2014}. También se observó que la agregación puede mejorar la precisión y reducir la incertidumbre en problemas de demanda, aunque disminuye la cantidad de observaciones disponibles \citep{Kourentzes2017}.

Este compromiso es especialmente relevante para el caso estudiado. La frecuencia semanal proporciona más puntos de entrenamiento, pero puede exhibir una variabilidad elevada. La frecuencia mensual reduce parte de ese ruido, aunque deja una muestra pequeña, mientras que el horizonte trimestral coincide mejor con la decisión financiera. Por lo tanto, la granularidad constituye una dimensión que debe evaluarse empíricamente.

El enfoque adoptado compara referencias simples, el pronóstico del CRM, modelos estadísticos, algoritmos basados en árboles y una red LSTM bajo un protocolo temporal común. La contribución no reside en proponer un algoritmo nuevo, sino en construir un proceso reproducible que transforma estados históricos del \textit{pipeline} en una proyección agregada, evita el uso de información futura y permite comparar familias de modelos sobre las mismas fechas y métricas.

\section{Objetivos}
\label{sec:objetivos}

El objetivo general del trabajo es desarrollar y evaluar un modelo de aprendizaje automático, integrado en un flujo de procesamiento de datos reproducible, para proyectar las contrataciones del trimestre siguiente a partir del estado mensual histórico de un CRM, con corrección temporal y un desempeño medido frente al pronóstico utilizado por la organización.

Dentro de los objetivos específicos se pueden mencionar los siguientes:

\begin{itemize}
    \item Diseñar e implementar un flujo reproducible para ingerir, validar, limpiar y consolidar los \textit{snapshots} históricos del CRM.
    \item Definir la variable objetivo de contrataciones y construir las variables predictoras con criterio \textit{point in time}, de modo de evitar duplicaciones y fuga de información.
    \item Analizar el efecto de las granularidades semanal, mensual y trimestral sobre la variabilidad de la serie y la precisión de los pronósticos.
    \item Comparar modelos de referencia, ARIMA, Prophet, \textit{Random Forest}, XGBoost y LSTM mediante una metodología común de validación temporal y métricas homogéneas.
    \item Seleccionar y ajustar el modelo de mejor desempeño y cuantificar su diferencia respecto del pronóstico del CRM.
    \item Generar proyecciones trimestrales agregadas, trazables e interpretables que puedan utilizarse como apoyo en la planificación financiera.
\end{itemize}

\section{Alcance del trabajo}
\label{sec:alcanceTrabajo}

El trabajo comprende el desarrollo de un prototipo funcional orientado a demostrar la factibilidad técnica y el valor práctico del enfoque. El alcance incluye la carga y validación de archivos históricos exportados desde el CRM, la consolidación de los \textit{snapshots}, la construcción de una tabla mensual de modelado, la definición de las contrataciones, la generación de variables, el entrenamiento y comparación de diferentes familias de modelos, el registro de experimentos y la exportación de proyecciones y métricas.

Quedaron fuera del alcance la conexión en tiempo real con el CRM, el desarrollo de una interfaz de usuario, el despliegue productivo y el modelado financiero posterior a la proyección. Tampoco se realizó una segmentación geográfica, debido a que la fuente disponible no contiene una dimensión de región.
